% ============================================================
% 01_intro.tex
% Expects \todo{} and \blocked{} macros defined in the preamble.
% Cross-references use labels: sec:related, sec:data, sec:divergence,
% sec:framework, sec:linking, sec:cost, sec:limitations, sec:conclusion.
% ============================================================

\section{Introduction}
\label{sec:intro}

Computational social science now compares machine-labelled text across platforms as a matter of course. \citet{overgaard2024} apply a BERT
classifier to Facebook and Twitter posts and read the difference in
scores as a difference in behaviour. \citet{wei2026} run the same
design across Douyin and TikTok. Studies in this mould publish a
cross-platform comparison, or a pooled multi-platform index, on the
assumption that the classifier behaves identically wherever it is
pointed. The assumption is rarely stated, still more rarely examined, and it is not an obviously safe one. \citet{nogara2025} report that Perspective API scores German comments as more toxic than equivalent English ones: the instrument behaves differently in different languages. In \citet{atreja2025}, LLM annotations shift under changes of prompt wording alone. And \citet{greene2025} find that the same political actors are read differently on different platforms. Classifiers are instruments, and their behaviour depends on the conditions in which they are used.

When the labels are machine-produced, nobody asks with a formal test whether the platform is part of the measurement. When a
classifier scores Reddit discussion of a budget more negative than
tweets about the same budget, the difference is read as a fact about opinion, when it could equally be a property of the instrument. Character
limits, anonymity, threading, moderation, and audience shape how the
same attitude is expressed, and therefore how it is scored.
\citet{sen2021} give this hazard a name, platform-affordance error,
in their total error framework for digital traces, but they provide no
formal test for it. Recent correction methods repair label error
against a human gold standard \citep{egami2023} and probe the
fragility of LLM annotation pipelines
\citep{baumann2025,halterman2025}, yet none of them models the
platform as a source of measurement error.

Survey methodology met a structurally identical problem decades ago
and built machinery for it. The same respondent answers differently by
telephone, on paper, and face to face; the field calls this a mode
effect and does not treat cross-mode differences as substantive until
measurement equivalence has been established
\citep{tourangeau2000,klausch2013,hox2015}, within a total-error
accounting of where distortions enter \citep{groves2010}. Our claim is that a platform is a mode. The analogy has limits: absent random assignment of users to platforms or within-user cross-posting, the design identifies platform-level non-equivalence, a bundle of selection and measurement that Section~\ref{sec:limitations} unpacks, with mode as the leading interpretation rather than an established fact. Affordances stand to posts as interview
settings stand to answers, and the psychometric toolkit built for modes and groups (generalizability theory \citep{cronbach1972,brennan2001}, measurement invariance testing \citep{meredith1993,putnick2016}, and score linking \citep{kolen2014}) transfers to machine-labelled constructs with the platforms as the groups. So far as we can find, no study has carried out that transfer: the toolkit has reached survey self-reports \citep{boyd2024},
annotator facets \citep{kennedy2020,sachdeva2022}, and single-platform
text scales \citep{pokropek2026}, but the platform itself has stayed
an unmodelled stratum.

The question is also practical. Since the closure of the major
research APIs, few teams can study one platform at depth; instead they
assemble patchworks of whatever platforms remain accessible, with
heterogeneous coverage and access terms
\citep{davidson2023,goanta2026}. Every such patchwork implicitly pools
or compares scores across platforms. If a unit of classifier score
does not mean the same thing on Reddit as on X, the pooled trend lines
and cross-platform contrasts now filling the literature carry an error term of unknown size.

This paper makes four contributions.
\begin{enumerate}
  \item We formalise the platform-as-mode question for machine-labelled
    constructs and give what we believe is the first adaptation of the
    psychometric equivalence toolkit to machine-classified policy
    sentiment with platform as the grouping facet.
  \item We specify a six-step protocol: same-event divergence testing;
    a fixed-facet generalizability study with a human-validated
    criterion; categorical measurement invariance and differential item
    functioning under effect-size criteria; multilevel decomposition of platform effects, with platform contrasts reported both with and without reaction latency; anchor-based score linking on pre-identified fiscal events,
    centred on the September 2022 mini-budget; and a simulation
    contrasting a naive pooled index with a measurement-adjusted one.
    All tests are replicated with a non-lexicon transformer, so only
    convergent non-invariance is read as a mode effect
    \citep{czarnek2022}, and invariance results are treated as
    diagnostic evidence rather than grounds to pool
    \citep{robitzsch2023}.
  \item We quantify the cost of assuming equivalence in a simulation with known ground truth: across 500 replications of one generating configuration fixed in advance, a pooled index that takes platform scores at face value misstates a known event shift by roughly an eighth, and because the platforms measure through different functions it reports the wrong sign of aggregate sentiment on nearly a quarter of days; a measurement-adjusted index recovers the shift with negligible bias and the correct sign on all but one day in sixty on average.
  \item We ground the protocol in a case study of the dissertation-stage instrument, a fine-tuned BERT classifier over 255,446 UK economic-policy comments from Twitter/X, Reddit, YouTube, and Quora, whose execution records supply the paper's motivating evidence; the full four-platform application, under fresh collection with the pre-specified design, is designated as the follow-up study.
\end{enumerate}

The paper proceeds as follows. Section~\ref{sec:related} locates the
problem between computational text measurement and survey mode-effect
research. Section~\ref{sec:data} describes the motivating corpus, the anchor events, and the validation design. Sections~\ref{sec:divergence} to~\ref{sec:linking} specify Steps 1 to 5, the protocol the registered application executes; Section~\ref{sec:cost} reports Step 6, the simulation that quantifies the cost of assuming equivalence.
Section~\ref{sec:limitations} states the limitations, Section~\ref{sec:ethics} sets out the ethics position, and Section~\ref{sec:conclusion} concludes.
